\question [20] O rompimento da barragem de contenção de uma mineradora em Mariana (MG) acarretou o derramamento de lama contendo resíduos poluentes no rio Doce. Esses resíduos foram gerados na obtenção de um minério composto pelo metal de menor raio atômico do grupo 8 da tabela de classificação periódica. A lama levou 16 dias para atingir o mar, situado a \SI{16}{\kilo\meter}  do local do acidente, deixando um rastro de destruição nesse percurso. Caso alcance o arquipélago de Abrolhos, os recifes de coral dessa região ficarão ameaçados. Com base nas informações apresentadas no texto, a velocidade média de deslocamento da lama, do local onde ocorreu o rompimento da barragem até atingir o mar, em \si[per-mode=symbol]{\kilo\meter\per\hour} corresponde a:   


		\begin{EnvUplevel}
			\begin{choices}
			
			\choice 3,8  
			\choice 2,1    
			\correctchoice 1,6       
			\choice 4,6      
			\choice 4,1 
				
			\end{choices}
		\end{EnvUplevel}
		
		\begin{EnvUplevel}
		\begin{solution}
		$$
		V_m=\frac{\Delta S}{\Delta T}=\frac{600}{24\times16}=1,56
		$$$$
		V_m=\SI[per-mode=fraction]{1,6}{\kilo\meter\per\hour}
		$$
		\end{solution}
		\end{EnvUplevel}


	
